\chapter{Conclusiones}
\label{cap:conclusiones}

Este trabajo desarrolló un sistema que identifica eventos sismo-volcánicos del
Complejo Volcánico Nevados de Chillán en registro continuo: detecta los
eventos con PhaseNet en la estación FRE, recorta cada uno según su energía,
genera su espectrograma y lo clasifica en VT, LP, TR o rechazo con una red
convolucional ajustada por transferencia de aprendizaje. El problema planteado
en el Capítulo~\ref{cap:introduccion} señalaba que faltaban tres cosas: evidencia
local de qué detector conviene usar, un recorte que se ajuste a la duración
real de los eventos y una medición de cuánto desempeño pierde el clasificador
al pasar de recortes de referencia a recortes automáticos. Las tres quedaron
medidas.

\section{Conclusiones por objetivo}

Sobre la detección (OE1), PhaseNet es el único de los tres detectores que
llega a la meta de 0,95 de sensibilidad volcánica sobre la traza continua
(0,952), la mantiene en tramos que no se usaron para elegir sus umbrales y la
supera en las ventanas de test (0,972). Lo logra gracias a la onda S y a un
umbral bajo, con un costo de unas 69 falsas alarmas por hora. EQTransformer
queda cerca con menos falsas alarmas, pero sus intervalos son demasiado cortos
para delimitar los eventos, y STA/LTA pierde casi todos los tremores al
recortar. También se mostró que la forma de contar una detección cambia el
resultado tanto como el detector; contar por intersección evita depender de
una tolerancia elegida a mano.

Sobre el recorte y los espectrogramas (OE2), el recorte por energía con corte
en el aviso siguiente delimita bien 0,864 de los volcánicos de la traza y casi
todos los VT y LP de test. El tremor es su límite: solo 0,65 de los TR quedan
bien recortados, porque el aviso llega tarde, porque PhaseNet avisa varias
veces dentro de un mismo tremor y porque algunos casi no sobresalen del ruido.
Con un buen recorte, la duración variable ayuda a reconocer el tremor; en
promedio, el recorte y la ventana fija de 60\,s rinden lo mismo.

Sobre el clasificador (OE3), VGG16, ResNet50 y EfficientNet-B0 alcanzan en
test un F1 macro de VT, LP y TR de 0,834 a 0,843, sin diferencias entre ellas
mayores que la variación entre folds. La búsqueda de la configuración y el
aumento de datos aportaron entre 0,02 y 0,03. La red liviana, EfficientNet-B0,
rinde igual que las grandes y clasifica más de tres veces más rápido que
VGG16. La confusión principal es entre TR y LP, y el clasificador manda a la
clase de rechazo solo una parte de las falsas alarmas del detector.

Sobre el sistema completo (OE4), en la traza continua el sistema detecta,
recorta y clasifica bien entre 0,646 y 0,704 de los volcánicos, y procesa las
9,74 horas en menos de dos minutos en un computador personal. Frente al
recorte de referencia, el sistema pierde entre 0,056 y 0,098 de sensibilidad
volcánica, casi todo en TR; cuando el recorte automático calza con su evento,
clasificar con él rinde lo mismo que con el de referencia. La pérdida en VT y
en la clase de rechazo frente a test no viene del recorte, y queda sin
explicar con los datos disponibles.

En conjunto, el sistema cumple el objetivo general. Su desempeño en registro
continuo queda por debajo del que logra cada etapa por separado, como ocurrió
en trabajos anteriores con una sola estación, y las mediciones dicen dónde se
pierde: en el recorte de los tremores, en la clasificación de VT y en las
falsas alarmas que el clasificador no rechaza.

\section{Recomendaciones y trabajo futuro}

Las recomendaciones salen de lo que los resultados mostraron como más débil.

Entrenar la clase de rechazo también con recortes que no corresponden a ningún
evento, o agregar una etapa que descarte falsas detecciones antes del
clasificador, como la que usó \textcite{soto2016llaima}. Es la vía más directa
para bajar las falsas alarmas que hoy llegan como volcánicas, y debe medirse
cuánto cuesta en sensibilidad.

Mejorar el recorte de los tremores sin usar información del catálogo, por
ejemplo juntando los avisos que caen dentro de un mismo tremor o usando la
curva de detección de EQTransformer para fijar el fin del evento. Cualquier
cambio debe compararse con el recorte actual usando el mismo IoU y la misma
traza.

Incorporar más estaciones. La asociación de avisos entre estaciones es la
forma habitual de descartar falsas alarmas, y combinar estaciones mejora la
clasificación en otros volcanes \parencite{curilem2016estaciones}.

Revisar con un analista los VT y los eventos de rechazo de la traza que el
clasificador confunde, para saber si la caída frente a test se debe a cómo se
etiquetaron o a la señal misma.

Ajustar PhaseNet y EQTransformer con datos volcánicos, para que reconozcan la
llegada gradual de los LP y los TR sin depender de umbrales tan bajos.

Evaluar el sistema con más horas de registro continuo y más eventos, para
estrechar los intervalos de confianza, y probarlo recibiendo señal en tiempo
real.
